% ============================================================
% paper/sections/06_cycles_closure.tex
% Cycles, Closure, and Persistence
% ============================================================

\section{Cycles and Closure}

\subsection{Cycles as a Condition for Persistence}

In the Ontology-of-Continua framework, persistence over time is not
assumed. A continuum may satisfy instantaneous admissibility conditions
while failing to persist due to the absence of recurrent structure.

For the enterprise continuum $K_{\mathrm{EA}}$, persistence requires
the existence of a minimal set of cycles $C_{\mathrm{EA}}$ that close
over the internal time scale $\tau_{\mathrm{EA}}$.

These cycles are structural necessities. They are not processes to be
optimized, improved, or redesigned within the scope of the model.

\subsection{Required Cycle Classes}

The Executable Theory Specification defines the following required
cycle classes:

\begin{itemize}
  \item \textbf{Governance cycles}: decision, enforcement, and feedback
        loops that maintain rule coherence.
  \item \textbf{Value cycles}: production, delivery, and consumption
        loops that sustain enterprise identity.
  \item \textbf{Reproduction cycles}: mechanisms that regenerate
        capabilities, roles, or structural components.
  \item \textbf{Financial cycles}: inflow, allocation, and outflow
        processes required for material persistence.
  \item \textbf{Learning cycles}: mechanisms that update internal models
        in response to observed discrepancies.
\end{itemize}

Each class represents a category of closure rather than a specific
implementation.

\subsection{Closure Rule}

Closure is defined structurally as the existence of a recurrent path in
the enterprise state graph. Formally, closure requires that admissible
trajectories form closed paths over $\tau_{\mathrm{EA}}$.

The ETS expresses this condition as a \emph{graph-closed path} or a
\emph{recurrent pattern over the internal time scale}. No assumptions
are made regarding efficiency, speed, or optimality of closure.

\subsection{Failure of Closure}

Failure of one or more required cycles does not necessarily imply
immediate termination. However, sustained absence of closure leads to
violation of stability thresholds.

In such cases, the continuum may transiently exist but cannot be
considered persistent. Diagnosis remains defined only until stability
or termination conditions are violated.

\subsection{Cycles Versus Processes}

It is important to distinguish cycles from operational processes.
Processes may exist without forming closed loops, and cycles may be
realized through heterogeneous processes.

The model does not prescribe how cycles are implemented. It asserts
only that their structural closure is required for persistence.

\subsection{Implications for Diagnosis}

The presence or absence of cycles affects:
\begin{itemize}
  \item stability threshold components,
  \item continuity $k_{\mathrm{EA}}$,
  \item admissibility of further diagnosis.
\end{itemize}

Cycles therefore serve as a bridge between instantaneous admissibility
and temporal persistence within the enterprise continuum.
